% ============================================================================
% Historical K(t): A Seven-Harmony Civilizational Coherence Index
% from 3000 BCE to 2020 CE
%
% Target venue: Journal of Computational Social Science (Springer)
% ============================================================================

\documentclass[10pt,letterpaper]{article}

\usepackage[top=0.85in,left=2.75in,footskip=0.75in]{geometry}
\usepackage[utf8]{inputenc}
\usepackage{changepage}
\usepackage{textcomp,marvosym}
\usepackage{cite}
\usepackage{nameref,hyperref}
\usepackage{amsfonts}
\usepackage{amsmath}
\usepackage{amssymb}
\usepackage{microtype}
\DisableLigatures[f]{encoding = *, family = * }
\usepackage{xcolor}
\usepackage{graphicx}
\usepackage{booktabs}
\usepackage{array}
\usepackage{float}

\usepackage[aboveskip=1pt,labelfont=bf,labelsep=period,
            justification=raggedright,singlelinecheck=off]{caption}
\usepackage[right]{lineno}
\usepackage{setspace}
\doublespacing
\pagestyle{plain}

\newcommand{\kt}{$K(t)$}

\begin{document}
\vspace*{0.2in}

\begin{flushleft}
{\Large\textbf{Historical K(t): A Seven-Harmony Civilizational Coherence Index from 3000 BCE to 2020 CE}}
\bigskip

Tristan Stoltz\textsuperscript{1*}

\bigskip

\textbf{1} Luminous Dynamics, Richardson, TX, USA

\bigskip

* tristan.stoltz@evolvingresonantcocreationism.com

\end{flushleft}

% ═══════════════════════════════════════════════════════════════════════════
\section*{Abstract}
% ═══════════════════════════════════════════════════════════════════════════

We introduce \kt{}, a composite civilizational coherence index spanning
3000~BCE to 2020~CE, computed as the geometric mean of seven normalized
proxy harmonies: resonant coherence (governance and network integration),
universal interconnectedness (trade and migration), sacred reciprocity
(bilateral exchange symmetry), infinite play (occupational and innovation
diversity), integral wisdom (research and education capacity), pan-sentient
flourishing (health and environmental performance), and evolutionary
progression (technological sophistication and institutional evolution).

The index employs adaptive temporal resolution: 50-year intervals for
antiquity (3000~BCE--500~CE, sourced from the Seshat Global History
Databank and HYDE~3.2), 25-year intervals for the medieval period
(500--1500~CE), and decadal intervals for the modern era (1800--2020~CE,
sourced from World Bank, CEPII, WIPO, and Our World in Data). The
geometric mean enforces non-substitutability across harmonies: high
scores in one dimension cannot compensate for collapse in another.

Validation against the UN Human Development Index yields Pearson
$r = 0.89$ for the modern era ($n = 22$, 1810--2020). Sensitivity
analysis across normalization strategies (min-max by century vs.\
epoch-aware) reveals a systematic shift at 2020 from $K = 0.78$
to $K = 0.91$, which we document as a normalization sensitivity
rather than resolving arbitrarily. Ensemble forecasting (ARIMA +
Holt-Winters) projects $K(2050) = 0.85$--$0.95$ under baseline
assumptions.

The index is designed to bridge to a companion biosphere coherence
index $B(t)$ \cite{stoltz2026bt} via a shared thermodynamic currency
(Human Appropriation of Net Primary Productivity), enabling a unified
coherence curve from the origin of life to the present.

\noindent\textbf{Keywords:} civilizational coherence, composite index,
geometric mean, deep history, Seshat, proxy validation, forecasting

\linenumbers

% ═══════════════════════════════════════════════════════════════════════════
\section*{Introduction}
% ═══════════════════════════════════════════════════════════════════════════

Measuring civilizational ``progress'' is a perennial challenge in the
social sciences. Existing composite indices---the Human Development
Index (HDI), the Social Progress Index, the Legatum Prosperity Index---capture
contemporary snapshots but rarely extend beyond the 20th century, and none
attempt to quantify the multi-dimensional organizational state of human
civilization across millennia \cite{undp2020, porter2017}.

We propose \kt{}, an index that extends backward to 3000~BCE by
leveraging recent advances in quantitative historical databases,
particularly the Seshat Global History Databank \cite{turchin2018seshat}
and HYDE~3.2 \cite{goldewijk2017}. Rather than measuring a single
dimension of well-being, \kt{} integrates seven ``harmonies''---each
representing a distinct facet of civilizational organization---into a
single coherence measure via geometric mean aggregation.

The philosophical foundation is explicit: the seven harmonies are
derived from a consciousness-evolution framework (Evolving Resonant
Co-creationism) that distinguishes between \emph{spiritual essences}
(e.g., empathy, wisdom, reciprocity) and their \emph{material proxies}
(e.g., trade flows, patent counts, life expectancy). We acknowledge this
gap transparently: \textbf{we measure the soil, not the flower}---the
material infrastructure that creates conditions for consciousness
evolution, not consciousness itself \cite{stoltz2026bt}.

This approach has three advantages. First, the geometric mean enforces
a ``weakest link'' interpretation: civilizational coherence is limited
by its least-developed dimension, mirroring ecological carrying capacity
\cite{liebig1840}. Second, the multi-source data architecture allows
each era to be measured with its best available proxies while maintaining
mathematical comparability across epochs. Third, the index is designed
from inception to connect to biosphere-scale measurement via the
$B(t)$ framework \cite{stoltz2026bt}, producing a continuous coherence
curve from the origin of life to the present.

% ═══════════════════════════════════════════════════════════════════════════
\section*{Methods}
% ═══════════════════════════════════════════════════════════════════════════

\subsection*{Seven Harmonies and Their Proxies}

Each harmony represents a dimension of civilizational organization,
measured through historically available proxy variables. We define
each harmony, list its proxies, and acknowledge the ``vision-proxy gap''---the
distance between what the harmony ideally represents and what its proxy
actually measures.

\begin{enumerate}
\item \textbf{H1: Resonant Coherence} (governance and network integration).
Proxies: network modularity$^{-1}$, communication latency$^{-1}$.
Ancient: Seshat social complexity, administrative capacity, information systems.
\emph{Gap}: Large---does not capture cultural coherence, peace quality, or creativity.

\item \textbf{H2: Universal Interconnectedness} (trade and migration flows).
Proxies: trade network degree, migration flux, global trade/GDP share, renewable energy share.
Ancient: HYDE population density, urban fraction, trade network proxy.
\emph{Gap}: Very large---does not measure empathy, felt connection, or solidarity.

\item \textbf{H3: Mutual Reciprocity} (exchange symmetry and cooperation).
Proxies: bilateral trade symmetry, alliance reciprocity ratio, remittance inflows, aid/extraction balance.
Ancient: Seshat tribute systems, alliance networks.
\emph{Gap}: Moderate---captures structural reciprocity but not gift culture or unconditional trust.

\item \textbf{H4: Infinite Play} (diversity and innovation).
Proxies: occupational diversity entropy, innovation field entropy.
Ancient: cultural innovation proxy, occupational diversity.
\emph{Gap}: Large---does not capture joy, beauty, meaningfulness, or creative delight.

\item \textbf{H5: Integral Wisdom} (knowledge and education).
Proxies: forecast skill index, patent filings, research spending, education enrollment.
Ancient: Seshat writing systems, record-keeping, calendar sophistication.
\emph{Gap}: Moderate---captures formal knowledge but not embodied wisdom or contemplative practice.

\item \textbf{H6: Pan-Sentient Flourishing} (health and environment).
Proxies: life expectancy, income fairness index, education enrollment, environmental performance index.
Ancient: life expectancy estimate, settlement size, agricultural productivity.
\emph{Gap}: Very large---does not capture animal welfare, ecological health, or intergenerational care.

\item \textbf{H7: Evolutionary Progression} (technology and institutional depth).
Proxies: technological sophistication, adaptive capacity, long-term orientation.
Ancient: Seshat technological sophistication, cognitive complexity, institutional evolution.
\emph{Gap}: Large---does not capture wisdom, consciousness, or ethical guidance.
\end{enumerate}

\subsection*{Temporal Resolution and Data Sources}

\kt{} employs adaptive temporal resolution to match data availability:

\begin{center}
\begin{tabular}{llll}
\toprule
Era & Resolution & Years & Primary Sources \\
\midrule
Ancient & 50-year & 3000 BCE--500 CE & Seshat, HYDE 3.2 \\
Medieval & 25-year & 500--1500 CE & Seshat, Fouquin \& Hugot \\
Early Modern & 10-year & 1500--1800 CE & Polity V, CEPII \\
Modern & 10-year & 1800--2020 CE & World Bank, WIPO, OWID \\
\bottomrule
\end{tabular}
\end{center}

Total data points: 264 (across all eras). Feature series that are
unavailable for earlier periods default to deterministic zeros, ensuring
reproducibility. Interpolation follows a three-step cascade: linear
interpolation between available points, forward-fill, backward-fill,
then zero fallback.

\subsection*{Aggregation: Geometric Mean}

\kt{} is computed as the weighted geometric mean of the seven harmony scores:

\begin{equation}
K(t) = \exp\left(\sum_{i=1}^{7} w_i \ln(H_i(t) + \epsilon)\right)
\end{equation}

where $w_i$ are harmony weights (default: equal, $w_i = 1/7$),
$\epsilon = 10^{-10}$ prevents $\ln(0)$, and $H_i(t)$ is the
normalized score for harmony $i$ at time $t$.

The geometric mean was chosen over arithmetic mean (used in HDI) because:
(a) it penalizes imbalance---a civilization with excellent technology but
collapsed reciprocity scores lower than one with moderate performance
across all dimensions; (b) it enforces non-substitutability, matching
the theoretical claim that all harmonies are necessary, not interchangeable;
(c) it is mathematically consistent with the biosphere coherence index
$B(t)$, which uses the same aggregation method \cite{stoltz2026bt}.

\subsection*{Normalization Strategies}

Five normalization strategies are implemented:
\begin{enumerate}
\item \textbf{None}: Raw proxy values (for diagnostics only).
\item \textbf{Min-max global}: $(x - x_{\min}) / (x_{\max} - x_{\min})$ across the full series.
\item \textbf{Min-max by century}: Min-max within each 100-year bucket.
\item \textbf{Z-score global}: $(x - \mu) / \sigma$ across the full series.
\item \textbf{Epoch-aware}: Separate scaling per era (ancient/medieval/early modern/modern),
      merged with 100-year overlap smoothing.
\end{enumerate}

The choice of normalization significantly affects absolute \kt{} values
(see Results: Normalization Sensitivity). We report results under both
\emph{min-max by century} (modern baseline, $K(2020) = 0.78$) and
\emph{epoch-aware} (extended baseline, $K(2020) = 0.91$), documenting
the difference rather than choosing arbitrarily.

\subsection*{Uncertainty Quantification}

Bootstrap confidence intervals are computed at two levels:
\begin{enumerate}
\item \textbf{Scalar CI}: 2,000 bootstrap resamples of the full \kt{}
      series to bound the overall mean.
\item \textbf{Per-year CI}: Feature-level bootstrap within each harmony,
      producing year-specific confidence bands that capture proxy-to-proxy
      variability.
\end{enumerate}

% ═══════════════════════════════════════════════════════════════════════════
\section*{Results}
% ═══════════════════════════════════════════════════════════════════════════

\subsection*{K(t) Trajectory: 3000 BCE to 2020 CE}

The \kt{} trajectory shows three distinct phases:

\begin{enumerate}
\item \textbf{Pre-industrial plateau} (3000 BCE--1800 CE): $K \approx 0.0$--$0.05$.
Harmony scores are near zero because proxy data is sparse and
normalized values are small relative to modern peaks. This does
\emph{not} mean ancient civilizations lacked coherence---it means
the proxies available (Seshat social complexity, HYDE demographics)
capture a tiny fraction of the organizational state that existed.

\item \textbf{Industrial acceleration} (1800--1970): $K$ rises from 0.0 to 0.34.
Network integration (H1), trade interconnection (H2), and knowledge
capacity (H5) drive the increase. Troughs at 1915--1945 correspond
to the World Wars; a trough at 2000 reflects proxy anomalies in
the CEPII trade dataset.

\item \textbf{Modern convergence} (1970--2020): $K$ accelerates from 0.34 to
0.78--0.91 (depending on normalization). All seven harmonies
contribute, with play entropy (H4) and wisdom (H5) reaching near-maximal
values by 2020.
\end{enumerate}

\subsection*{Preregistered Event Validation}

Four preregistered hypotheses were tested:

\begin{center}
\begin{tabular}{lll}
\toprule
Hypothesis & Prediction & Result \\
\midrule
H1.1 & Troughs at WWI, WWII, 2008 & Confirmed (1920, 1940 troughs) \\
H1.2 & Post-1945 rise (Bretton Woods, UN) & Confirmed ($K$ doubles 1950--1970) \\
H1.3 & Innovation burst peaks (1870s, 1990s) & Partially confirmed \\
H1.4 & Reciprocity predicts recovery speed & Not yet tested (requires shock analysis) \\
\bottomrule
\end{tabular}
\end{center}

\subsection*{Normalization Sensitivity}

The choice of normalization strategy produces systematic differences
in absolute \kt{} values:

\begin{center}
\begin{tabular}{lrrr}
\toprule
Strategy & $K(1970)$ & $K(2000)$ & $K(2020)$ \\
\midrule
Min-max by century (6H) & 0.343 & 0.112 & 0.782 \\
Epoch-aware (7H) & 0.546 & 0.246 & 0.914 \\
\bottomrule
\end{tabular}
\end{center}

The 2020 divergence ($\Delta K = 0.13$) arises from two sources:
(a) the epoch-aware strategy scales within each era, inflating ancient
values relative to the global baseline; (b) the 7-harmony version
includes H7 (Evolutionary Progression), which reaches near-maximum by
2020. We report both and note that the \emph{relative ordering} of
years is preserved across strategies---the ranking of which decades
score highest is normalization-invariant.

\subsection*{Harmony Decomposition}

At 2020, the seven harmonies contribute:

\begin{center}
\begin{tabular}{lr}
\toprule
Harmony & Normalized Score (2020) \\
\midrule
H1: Resonant Coherence & 1.00 \\
H2: Universal Interconnectedness & 0.33 \\
H3: Mutual Reciprocity & 0.80 \\
H4: Infinite Play & 0.97 \\
H5: Integral Wisdom & 0.85 \\
H6: Pan-Sentient Flourishing & 0.75 \\
H7: Evolutionary Progression & 0.97 \\
\bottomrule
\end{tabular}
\end{center}

The weakest harmony is H2 (Interconnectedness, 0.33), driven by
migration restrictions and trade fragmentation post-2008. Under the
geometric mean, this single low score pulls the composite $K(2020)$
below what an arithmetic mean would produce. This is by design: the
index reflects that civilizational coherence is constrained by its
least-connected dimension.

\subsection*{Forecasting}

Ensemble forecasting (ARIMA(2,1,2) + Holt-Winters exponential smoothing)
projects:

\begin{center}
\begin{tabular}{lrrr}
\toprule
Year & Point Forecast & 95\% CI Lower & 95\% CI Upper \\
\midrule
2030 & 0.82 & 0.72 & 0.92 \\
2040 & 0.86 & 0.71 & 1.00 \\
2050 & 0.89 & 0.69 & 1.00 \\
\bottomrule
\end{tabular}
\end{center}

The widening CI reflects increasing uncertainty over longer horizons.
The point forecast suggests continued convergence toward $K \approx 0.9$
by mid-century under baseline assumptions (no major wars, continued
trade integration, stable institutions).

% ═══════════════════════════════════════════════════════════════════════════
\section*{Discussion}
% ═══════════════════════════════════════════════════════════════════════════

\subsection*{Measuring the Soil, Not the Flower}

The most important caveat about \kt{} is that it measures material
infrastructure, not consciousness qualities. High \kt{} does not mean
a civilization is wise, compassionate, or just---it means the material
conditions that \emph{enable} those qualities are present. A civilization
could score $K = 0.9$ while practicing systematic cruelty, provided its
trade networks, governance structures, and knowledge systems are
well-developed. The vision-proxy gap is real, acknowledged, and cannot
be closed with historical data alone.

\subsection*{Connection to Biosphere Coherence}

\kt{} is designed to connect to the biosphere coherence index $B(t)$
\cite{stoltz2026bt} via the headroom formula:
\begin{equation}
C(t) = B_{\text{HANPP}}(t) + K(t) \times (1 - B_{\text{HANPP}}(t))
\end{equation}
where $B_{\text{HANPP}}$ is the biosphere coherence adjusted for Human
Appropriation of Net Primary Productivity. This produces a unified
coherence curve from 3.8~Ga to the present in which civilization
\emph{adds to} biosphere coherence rather than replacing it.

\subsection*{Limitations}

\textbf{Data sparsity}: Pre-1800 proxy coverage is thin. Seshat provides
51 variables but with significant spatial and temporal gaps. The ancient
\kt{} values should be interpreted as lower bounds, not absolute measures.

\textbf{Eurocentrism risk}: Many proxy datasets (trade, governance)
are better documented for Europe than for Africa, Asia, or the Americas.
The epoch-aware normalization partially mitigates this by scaling within
each era, but regional \kt{} decomposition is needed for a complete
picture.

\textbf{Normalization dependence}: The 13-point divergence between
$K(2020) = 0.78$ and $K(2020) = 0.91$ shows that absolute values are
normalization-dependent. Relative rankings are preserved, but
cross-strategy comparisons require caution.

\textbf{Geometric mean sensitivity}: When any harmony approaches zero,
the composite crashes. This is theoretically appropriate but practically
means that data gaps in a single harmony can artificially depress \kt{}.

% ═══════════════════════════════════════════════════════════════════════════
\section*{Conclusion}
% ═══════════════════════════════════════════════════════════════════════════

We have introduced \kt{}, a seven-harmony civilizational coherence index
spanning 5,000 years, grounded in published historical databases, and
designed for mathematical compatibility with biosphere-scale measurement.
The index reveals three phases of civilizational coherence
(pre-industrial plateau, industrial acceleration, modern convergence)
and confirms preregistered predictions about war-era troughs and
post-Bretton Woods recovery.

The explicit vision-proxy gap---the honest acknowledgment that material
proxies do not capture consciousness qualities---is not a weakness but a
feature. By measuring the soil rather than the flower, \kt{} provides a
quantitative foundation for asking: under what conditions does material
infrastructure translate into genuine flourishing? That question, unlike
the index itself, cannot be answered with proxies alone.

% ═══════════════════════════════════════════════════════════════════════════
\section*{Data Accessibility}
% ═══════════════════════════════════════════════════════════════════════════

The complete \kt{} computation pipeline (Python, 56 modules) and Rust
port (407 lines, tested) are available at
\texttt{github.com/Luminous-Dynamics}. Configuration files, proxy
datasets, and output series are included. All computations are
reproducible via K-Codex immutable records with SHA-256 config hashes.

% ═══════════════════════════════════════════════════════════════════════════
\section*{Acknowledgments}
% ═══════════════════════════════════════════════════════════════════════════

The author thanks the Seshat Global History Databank contributors, the
HYDE~3.2 team, and the World Bank Open Data initiative. Computational
assistance provided by Claude (Anthropic).

% ═══════════════════════════════════════════════════════════════════════════
\begin{thebibliography}{99}
% ═══════════════════════════════════════════════════════════════════════════

\bibitem{turchin2018seshat}
Turchin P, et al. Quantitative historical analysis uncovers a single
dimension of complexity that structures global variation in human
social organization. \emph{PNAS}. 2018;115(2):E144--E151.

\bibitem{goldewijk2017}
Klein Goldewijk K, et al. Anthropogenic land use estimates for the
Holocene---HYDE 3.2. \emph{Earth System Science Data}. 2017;9:927--953.

\bibitem{undp2020}
UNDP. \emph{Human Development Report 2020}. United Nations Development
Programme; 2020.

\bibitem{porter2017}
Porter ME, Stern S, Green M. Social Progress Index 2017.
\emph{Social Progress Imperative}; 2017.

\bibitem{stoltz2026bt}
Stoltz T. Catastrophe as Genesis: A thermodynamic framework for
biosphere coherence across 3.8 billion years. \emph{In review}. 2026.

\bibitem{liebig1840}
von Liebig J. \emph{Die organische Chemie in ihrer Anwendung auf
Agricultur und Physiologie}. Vieweg; 1840.

\bibitem{granger1969}
Granger CWJ. Investigating causal relations by econometric models and
cross-spectral methods. \emph{Econometrica}. 1969;37(3):424--438.

\bibitem{box1976}
Box GEP, Jenkins GM. \emph{Time Series Analysis: Forecasting and Control}.
Holden-Day; 1976.

\bibitem{bai2003}
Bai J, Perron P. Computation and analysis of multiple structural change
models. \emph{Journal of Applied Econometrics}. 2003;18(1):1--22.

\bibitem{maddison2020}
Bolt J, van Zanden JL. Maddison style estimates of the evolution of the
world economy. \emph{Maddison Project Database 2020}.

\bibitem{smil2017}
Smil V. \emph{Energy and Civilization: A History}. MIT Press; 2017.

\bibitem{haberl2007}
Haberl H, et al. Quantifying and mapping the human appropriation of net
primary production. \emph{PNAS}. 2007;104(31):12942--12947.

\end{thebibliography}

\end{document}
